\subsection{Inner Product Spaces, Euclidean Spaces}

\begin{tcolorbox}[title=Problem 1, breakable]
    Show that if $x$ and $y$ are vectors in an inner product
        space such that $||x||^2 = ||y||^2 = (x|y)$ then $x = y$.
    What does this mean for the inner product space of Example 5.1.3?
\end{tcolorbox}

\begin{proof}
    Suppose $x$ and $y$ are vectors in an inner product space.
    We have 
    \[(x|x) = (y|y) = (x|y).\]
    Now
    \begin{align*}
        \|x - y\|^2 &= (x - y \mid x - y) \\
        &= (x \mid x) - (x \mid y) - (y \mid x) + (y \mid y).
    \end{align*}
    Since $(x|y) = (y|x)$
    \[\|x - y\|^2 = (x \mid x) - (x \mid y) - (x \mid y) + (y \mid y) = 2(x \mid y) - 2(x \mid y) = 0.\]
    Therefore $x - y = 0$, so $x = y$.
\end{proof}


\begin{tcolorbox}[title=Problem 2, breakable]
    (\emph{Bessel's Inequality}) In an inner product space, if $x_1, \ldots, x_n$
    are pairwise orthogonal unit vectors $(||x_i|| = 1$ for all $i$,
    and $x_i \perp x_j$ when $i \ne j$), then 
    \[\sum_{i = 1}^{n} |(x | x_i)|^2 \le ||x||^2.\]
    (For $n = 1$ this is essentially the Cauchy-Schwarz inequality.)
    {Hint: Define $c_i = (x|x_i), y = c_1 x_1 + \cdots + c_n x_n$ and 
    $z = x - y$. Show that $y \perp z$ and apply the Pythagorean formula to 
    $x = y + z$ to conclude that $||x||^2 \ge ||y||^2$.}
\end{tcolorbox}

\begin{proof}
    Suppose in an inner product space, $x_1, \ldots, x_n$ are pairwise orthogonal unit vectors.
    Define $c_i = (x|x_i), y = \sum c_i x_i$ and 
        $z = x - y$.
    So $y = \sum (x|x_i) x_i$ and $z = x - \sum (x|x_i) x_i$
    We have 
    \[(y|z) = \left(\sum_i c_i x_i \,\Big|\, x - \sum_i c_i x_i\right) = \left(\sum_i c_i x_i \,\Big|\, x\right) - \left(\sum_i c_i x_i \,\Big|\, \sum_i c_i x_i\right).\]
    Now,
    \[\left(\sum_i c_i x_i \,\Big|\, \sum_i c_i x_i\right) = \sum_i c_i^2 \left(x_i|x_i\right) = \sum_i c_i^2.\]
    Then 
    \[\left(\sum_i c_i x_i \,\Big|\, x\right) - \left(\sum_i c_i x_i \,\Big|\, \sum_i c_i x_i\right) = \sum_i c_i (x_i | x) - \sum_i c_i^2 = \sum_i c_i^2 - \sum_i c_i^2 = 0.\]
    Thus $y \perp z$. So 
    \[\|x\|^2 = \|y + z\|^2 = \|y\|^2 + \|z\|^2 \geq \|y\|^2 = \sum_i c_i^2\].
\end{proof}

\begin{tcolorbox}[title=Problem 3, breakable]
    For nonzero vectors $x, y$ in an inner product space 
    $||x + y|| = ||x|| + ||y||$ if and only if $y = cx$ with $c > 0$.
\end{tcolorbox}

\begin{proof}
    Suppose $\|x + y\| = \|x\| + \|y\|$.
    Squaring both sides we have $\|x + y\|^2 = \|x\|^2 + \|y\|^2 + 2(x|y)$.
    Also $(\|x\| + \|y\|)^2 = \|x\|^2 + \|y\|^2 + 2\|x\|\|y\|$, so $(x|y) = \|x\|\|y\|$.
    By the Cauchy-Schwarz equality condition, this holds if and only if $y = cx$ for some scalar $c$,
    and since $(x|y) = c\|x\|^2 \geq 0$ we need $c > 0$.

    Conversely, suppose $y = cx$ with $c > 0$.
    Then $\|x + y\| = \|x + cx\| = (1+c)\|x\|$ and $\|x\| + \|y\| = \|x\| + c\|x\| = (1+c)\|x\|$.
\end{proof}


\begin{tcolorbox}[title=Problem 4, breakable]
    Vectors $x, y$ in a (real) inner product space are orthogonal if and only if 
    $||x + y|| = ||x - y||$.
\end{tcolorbox}

\begin{proof}
    Let $x, y$ be vectors in a real inner product space. Then 
    \begin{align*}
        \|x + y\| = \|x - y\| &\iff \|x + y\|^2 = \|x - y\|^2 \\
        &\iff \|x\|^2 + 2(x|y) + \|y\|^2 = \|x\|^2 - 2(x|y) + \|y\|^2 \\
        &\iff 4(x|y) = 0 \iff (x|y) = 0.
    \end{align*}
\end{proof}

\begin{tcolorbox}[title=Problem 5, breakable]
    In an inner product space, let $x_1, \ldots, x_n$ be vectors that are 
    `pairwise orthogonal', that is, $(x_i|x_j) = 0$ whenever $i \ne j$.
    \begin{enumerate}
        \item Prove that $|| \sum_{i = 1}^{n} x_i ||^2 = \sum_{i = 1}^{n} ||x_i||^2$.
        \item Deduce an alternative proof of Theorem 5.1.12.
    \end{enumerate}
\end{tcolorbox}

\begin{proof}
    Expanding the lhs we reveal 
    \[\sum_{i = 1}^{n} ||x_i||^2 + \text{ cross terms.}\]
    Since $(x_i|x_j) = 0$ where $i \ne j$ the cross terms are all $0$.
\end{proof}

\begin{proof}
    Suppose $x_1, \ldots, x_n$ are pairwise orthogonal nonzero vectors in an inner product space.
    Suppose $\sum_{i=1}^{n} \lambda_i x_i = 0$. By (1),
    \[\sum_{i = 1}^{n} ||\lambda_i x_i||^2 = \left\|\sum_{i = 1}^{n} \lambda_i x_i\right\|^2 = 0.\]
    Since each $||x_i|| \ne 0$, we have $\lambda_i = 0$ for all $i$, so the vectors are linearly independent.
\end{proof}


\begin{tcolorbox}[title=Problem 6, breakable]
    If $r_1, \ldots r_n$ are real numbers $>0$, then the formula 
    \[[x|y] = \sum_{i = 1}^{n} r_i a_i b_i,\]
    for $x = (a_1, \ldots, a_n), y = (b_1, \ldots, b_n)$, defines an inner 
    product on $\mathbb{R}^n$.
\end{tcolorbox}

\begin{proof}
    Let $x = (a_1, \ldots, a_n), y = (b_1, \ldots, b_n), z = (c_1, \ldots, c_n) \in \mathbb{R}^n$.
    Let $c$ be a scalar.
    Suppose $x \ne 0$ then 
    \[[x|x] = \sum_{i = 1}^{n} r_i a_i a_i \ne 0\]
    since some $a_i \ne 0$ and $r_i \ne 0$.
    Clearly 
    \[[x|y] = \sum_{i = 1}^{n} r_i a_i b_i = \sum_{i = 1}^{n} r_i b_i a_i = [y|x],\]
    since $\mathbb{R}$ is a field.
    Also 
    \[[x + y|z] = \sum_{i=1}^{n} r_i (a_i + b_i) c_i = \sum_{i=1}^{n} r_i a_i c_i + \sum_{i=1}^{n} r_i b_i c_i = [x|z] + [y|z],\]
    and 
    \[[cx|y] = \sum_{i=1}^{n} r_i (c a_i) b_i = c \sum_{i=1}^{n} r_i a_i b_i = c[x|y].\]
\end{proof}

\begin{tcolorbox}[title=Problem 7, breakable]
    In the inner product space of Example 5.1.3, if $x$ and $y$ are continuously 
    differentiable then 
    \[(x|y') + (x'|y) = x(b)y(b) - x(a)y(a).\]
\end{tcolorbox}

\begin{proof}
    We have 
    \[(x|y') = \int_{a}^{b} x(t) y'(t) \, dt \text{ and } (x'|y) = \int_{a}^{b} x'(t) y(t) \, dt,\]
    so 
    \[(x|y') + (x'|y) = \int_{a}^{b} x(t) y'(t) + x'(t) y(t) \, dt = \int_{a}^{b} \frac{d}{dt}[x(t)y(t)] \, dt = \Big[x(t)y(t)\Big]_{a}^{b} = x(b)y(b) - x(a)y(a). \qedhere\]
\end{proof}


\begin{tcolorbox}[title=Problem 8, breakable]
    Prove that, in a (real) inner product space, $||x|| = ||y||$ if and only if 
    $x + y$ and $x - y$ are orthogonal. Geometric Interpretations?
\end{tcolorbox}

\begin{proof}
    We have 
    \[\|x\| = \|y\| \iff (x|x) = (y|y) \iff (x|x) - (y|y) = 0 \iff (x+y|x-y) = 0 \iff x+y \perp x-y,\]
    where the third equivalence uses $(x+y|x-y) = (x|x) + (y|x) - (x|y) - (y|y) = (x|x) - (y|y)$.
\end{proof}

\begin{tcolorbox}[title=Problem 9, breakable]
    If $T$ is a linear mapping in a (real) inner product space, then 
    \[(Tx|Ty) = \frac{1}{4}\{||T(x + y)||^2 - ||T(x - y)||^2\}\]
    for all $x, y$ by $x + y, x - y$.
\end{tcolorbox}

\begin{proof}
    Suppose $T$ is a linear mapping in a (real) inner product space.
    Notice 
    \[||T(x + y)||^2 = ||Tx + Ty||^2 = (Tx + Ty | Tx + Ty) = (Tx|Tx) + 2(Tx|Ty) + (Ty|Ty),\]
    and similarly,
    \[||T(x - y)||^2 = (Tx - Ty | Tx - Ty) = (Tx|Tx) - 2(Tx|Ty) + (Ty|Ty).\]
    Subtracting and dividing by $4$ gives
    \[\frac{1}{4}\{||T(x+y)||^2 - ||T(x-y)||^2\} = \frac{1}{4} \cdot 4(Tx|Ty) = (Tx|Ty). \qedhere\]
\end{proof}

\begin{tcolorbox}[title=Problem 10, breakable]
    For subsets $A, B$ of an inner product space, (i) $A \perp B \iff A \subset B^\perp$,
    and (ii) $(A + B)^\perp = (A \cup B)^\perp = A^\perp \cap B^\perp$. (iii) In general,
    $(A \cap B)^\perp \ne A^\perp + B^\perp$.
\end{tcolorbox}

\begin{proof}
    For all $x \in A$ and $y \in B$,
    \[A \perp B \iff (x|y) = 0 \text{ for all } x \in A, y \in B \iff x \in B^\perp \text{ for all } x \in A \iff A \subset B^\perp.\]
\end{proof}

\begin{proof}
    Since $A, B \subset A \cup B \subset A + B$ by Theorem 5.1.14 (iii) we have 
    \[(A + B)^\perp \subset (A \cup B)^\perp \subset A^\perp \cap B^\perp.\]
    Now let $x \in A^\perp \cap B^\perp$. So for all $y \in A$, $(x|y) = 0$, and for all $y \in B$, $(x|y) = 0$.
    Thus if $y \in A \cup B$, $(x|y) = 0$ thus $x \in (A \cup B)^\perp$.
    Suppose $c = a + b \in A + B$. Then $(x|a + b) = (x|a) + (x|b) = 0 + 0 = 0$
        thus $c \in (A + B)^\perp$ completing the proof.
\end{proof}

\textbf{Solution (iii): }
In $\mathbb{R}^2$ consider $A = \{e_1, e_2\}$ and $B = \{e_1, e_1 + e_2\}$.
Then $A \cap B = \{e_1\}$, so
\[A^\perp + B^\perp = \{0\} + \{0\} = \{0\}\]
but 
\[(A \cap B)^\perp = \{e_1\}^\perp = \operatorname{span}\{e_2\} \ne \{0\}.\]

\subsection{Duality in Inner Product Spaces}

\begin{tcolorbox}[title=Problem 1, breakable]
Let $E$ be the inner product space of 5.1.3 (the space of continuous real-valued functions on a closed interval $[a, b]$, with inner products defined by means of the integral).
\begin{enumerate}[(i)]
    \item Fix a point $c \in [a, b]$ and let $f$ be the linear form on $E$ defined by $f(x) = x(c)$ for all $x \in E$ ($f$ is `pointwise evaluation at $c$'). Prove that there does \textit{not} exist a function $y \in E$ such that $f = y'$ in the sense of Definition 5.2.1. \{Hint: Note that $f \neq 0$. Show that if $y \in E$, $y \neq 0$, then there exists $x \in E$ such that $x(c) = 0$ but $\int_a^b x(t)y(t)\,dt > 0$.\}

    \item Let $f$ be the linear form on $E$ defined by
    \[
        f(x) = \int_a^b x(t)\,dt \quad \text{for } x \in E.
    \]
    True or false (explain): There exists $y \in E$ such that $f = y'$.

    \item Fix a point $c$, $a < c < b$, and let $f$ be the linear form on $E$ defined by
    \[
        f(x) = \int_a^c x(t)\,dt \quad \text{for } x \in E.
    \]
    Same question as in (ii). (Different answer!)
\end{enumerate}
\end{tcolorbox}

\begin{proof}
Consider $x \in E$ defined by $x(t) = 1$, then $f(x) = x(c) = 1$.
Thus $f \neq 0$. Suppose, for contradiction, there exists $y \in E$ such
that $f = y'$. That is, for all $x \in E$, we have $f(x) = y'(x) =
(y \mid x)$. So,
\[
    f(x) = x(c) = y'(x) = (y \mid x) = \int_a^b x(t)y(t)\,dt.
\]
Consider the function $x(t) = (t-c)^2 y(t)$. We have $x(c) = 0 = f(x)$.
But
\[
    \int_a^b x(t)y(t)\,dt = \int_a^b (t-c)^2 y(t)^2\,dt.
\]
Since $y \neq 0$ we must have $\int_a^b x(t)y(t)\,dt > 0$. Thus
$(x \mid y) \neq f(x)$, which is a contradiction.
\end{proof}

\begin{proof}
Let $y = 1 \in E$. Then for all $x \in E$,
\[
    (y \mid x) = \int_a^b x(t)y(t)\,dt = \int_a^b x(t)\,dt = f(x).
\]
Thus $f = y'$ with $y = 1$.
\end{proof}


\begin{proof}
Suppose, for contradiction, there exists $y \in E$ such that $f = y'$ so
\[
    \int_a^c x(t)\,dt = \int_a^b x(t)y(t)\,dt \quad \text{for all } x \in E.
\]
Consider the function
\[
    x(t) = \begin{cases} 0 & t \in [a, c] \\ (t-c)y(t) & t \in [c, b]. \end{cases}
\]
Then
$f(x) = 0$, so we require
\[
    \int_c^b (t-c)y(t)^2\,dt = 0.
\]
Since $(t-c) > 0$ on $(c,b]$ so $y = 0$ on $[c,b]$, thus
$y(c) = 0$. Now consider
\[
    x(t) = \begin{cases} (t-c)^2 & t \in [a, c] \\ 0 & t \in [c, b]. \end{cases}
\]
Then
\[
    \int_a^c (t-c)^2\,dt = \int_a^c (t-c)^2 y(t)\,dt,
\]
so $\int_a^c (t-c)^2(1 - y(t))\,dt = 0$. Since $(t-c)^2 > 0$ on $[a,c)$,
we have $y = 1$ on $[a,c]$ thus $y(c) = 1$. This contradicts
$y = 0$.
\end{proof}

\begin{tcolorbox}[title=Problem 2, breakable]
Let $V$ be the vector space of all functions $x : \mathbf{P} \to \mathbf{R}$ 
with the pointwise linear operations (cf.\ 1.3.7) 
and let $E$ be the set of all functions $x \in V$ whose `support'
\[
    \{n \in \mathbf{P} : x(n) \neq 0\}
\]
is a finite subset of $\mathbf{P}$. Prove:
\begin{enumerate}[(i)]
    \item $E$ is a linear subspace of $V$;
    \item the formula
    \[
        (x \mid y) = \sum_n x(n)y(n)
    \]
    defines an inner product on $E$ ($n$ runs over $\mathbf{P}$, but the sum has at most finitely many nonzero terms).
    \item The formula
    \[
        f(x) = \sum_n x(n)
    \]
    defines a linear form on $E$.
    \item There does not exist a vector $y \in E$ such that $f = y'$.
\end{enumerate}
\end{tcolorbox}

\begin{proof}
    Let $f, g \in E$ and let $c$ be a scalar.
    We want to show that $cf$ and $f + g$ have supports with finite cardinalities.
    If $c = 0$, then $cf$ is the zero function, whose support is empty, thus finite.
    If $c \ne 0$, then $(cf)(n) \ne 0 \iff f(n) \ne 0$, so $cf$ has the same support as $f$, which is finite by assumption.
    Notice that if $(f + g)(n) \ne 0$, then at least one of $f(n)$ or $g(n)$ is nonzero.
    Therefore
    \[
        \operatorname{supp}(f + g) \subseteq \operatorname{supp}(f) \cup \operatorname{supp}(g).
    \]
    Since $\operatorname{supp}(f)$ and $\operatorname{supp}(g)$ are both finite by assumption, their union is finite, and so $\operatorname{supp}(f + g)$ is finite.
\end{proof}

\begin{proof}
    Consider $f, g, t \in E$.
    \[(f \mid f) = \sum_n f(n)f(n).\]
    Since $f(n) \ne 0$ for $n$ in the support we have $f(n)f(n) = (f(n))^2 > 0$,
    and all other terms are $0$, so $(f \mid f) \ge 0$ with equality iff $f = 0$.
    Then 
    \[(f \mid g) = \sum_n f(n) g(n) = \sum_n g(n) f(n) = (g \mid f).\]
    Then 
    \[(f + g \mid t) = \sum_n [f(n) + g(n)]t(n) = \sum_n f(n)t(n) + g(n)t(n) = \sum_n f(n)t(n) + \sum_n g(n)t(n) = (f\mid t) + (g \mid t).\]
    Finally
    \[(cf \mid g) = \sum_n ((cf)(n))g(n) = c \sum_n f(n)g(n) = c(f \mid g).\]
\end{proof}

\begin{proof}
    We must show $f$ is additive and homogeneous.
    Let $x, y \in E$ and $c \in \mathbb{R}$, then
    \[f(x + y) = \sum_n (x + y)(n) = \sum_n [x(n) + y(n)] = \sum_n x(n) + \sum_n y(n) = f(x) + f(y).\]
    Then
    \[f(cx) = \sum_n (cx)(n) = c \sum_n x(n) = cf(x).\]
    Thus $f$ is a linear form on $E$.
\end{proof}

\begin{proof}
    For contradiction, suppose there exists a vector $y \in E$ such that $f = y'$, for all $x \in E$.
    Since $y \in E$, its support is finite. Let $N = \max\operatorname{supp}(y)$ (or $N = 0$ if $y = 0$).
    Define $x \in E$ as the function with $x(N+1) = 1$ and $x(n) = 0$ otherwise.
    Then $x \in E$ and
    \[f(x) = \sum_n x(n) = 1,\]
    but
    \[(x \mid y) = \sum_n x(n)y(n) = y(N+1) = 0,\]
    since $N+1 > N \ge \max\operatorname{supp}(y)$.
    This contradicts $f(x) = (x \mid y)$, so no such $y \in E$ exists.
\end{proof}


\begin{tcolorbox}[title=Problem 4, breakable]
If $E$ is any inner product space (not necessarily finite-dimensional), then $M = M^{\perp\perp}$ for every finite-dimensional linear subspace $M$. \{Hint: The proof of Corollary 5.2.11 is not applicable here. The problem is to show that $M^{\perp\perp} \subset M$. If $x \in M^{\perp\perp}$ and $x = y + z$ as in the proof of Theorem 5.2.7, then $z = x - y \in M^{\perp\perp}$, so $z \perp z$.\}
\end{tcolorbox}

\begin{proof}
    Suppose $E$ is any inner product space. Let $M$ be a finite-dimensional linear subspace of $E$.
    We know by Theorem 5.1.14, $M \subset M^{\perp \perp}$.
    Furthermore, let $x \in M^{\perp \perp}$.
    Since $M$ is a finite-dimensional subspace of $E$, by Theorem 5.2.7 we can write $x = y + z$
        for some $y \in M$ and $z \in M^{\perp}$.
    So $z = x - y$ where $x \in M^{\perp \perp}$ and $y \in M \subset M^{\perp\perp}$,
        thus $z \in M^{\perp\perp}$.
    By Theorem 5.1.14 (v), $z \in M^{\perp} \cap M^{\perp\perp} \subset \{0\}$ so 
    \[z = 0 \implies z = 0 \implies x = y \in M.\]
    Thus $M^{\perp\perp} \subset M$.
\end{proof}

\begin{tcolorbox}[title=Problem 5, breakable]
It can be shown that in the inner product space of continuous functions on $[a, b]$ (Example 5.1.3), the polynomial functions form a linear subspace $M$ such that $M^\perp = \{0\}$ (the proof is not easy!), therefore $M \neq M^{\perp\perp}$.
\end{tcolorbox}

\begin{tcolorbox}[title=Problem 6, breakable]
If $M$ and $N$ are linear subspaces of a Euclidean space, then
\[
    (M + N)^\perp = M^\perp \cap N^\perp \quad \text{and} \quad (M \cap N)^\perp = M^\perp + N^\perp.
\]
\{Hint: Corollary 5.2.11.\}
\end{tcolorbox}

\begin{proof}
    Since $M, N$ are subsets of an inner product space,
    by Section 5.1 Exercise 10 we have $(M + N)^\perp = M^\perp \cap N^\perp$.
\end{proof}
\begin{proof}
    We have 
    \begin{align*}
    (M \cap N)^\perp &= (M^{\perp \perp} \cap N^{\perp \perp})^\perp && \text{Corollary 5.2.11} \\
                     &= ((M^\perp + N^\perp)^\perp)^\perp && \text{Part 1} \\
                     &= M^\perp + N^\perp && \text{Corollary 5.2.11}
    \end{align*}
\end{proof}


\begin{tcolorbox}[title=Problem 8, breakable]
\begin{enumerate}[(i)]
    \item Let $E$ be an inner product space and let $x_1, \ldots, x_n, x$ be vectors such that $x$ is not a linear combination of $x_1, \ldots, x_n$. Prove that there exists a vector $y$ such that
    \[
        [\{x_1, \ldots, x_n, x\}] = [\{x_1, \ldots, x_n, y\}]
    \]
    and such that $y$ is orthogonal to every $x_i$. \{Hint: Exercise 3, part (ii).\}

    \item If $x_1, \ldots, x_n$ are linearly independent vectors in $E$, use the idea of (i) to construct orthonormal vectors $y_1, \ldots, y_n$ such that $[\{x_1, \ldots, x_k\}] = [\{y_1, \ldots, y_k\}]$ for $k = 1, \ldots, n$. (See Exercise 12 for a more explicit construction.)
\end{enumerate}
\end{tcolorbox}

\begin{proof}
    Let $M = [\{x_1, \ldots, x_n\}]$, $N = [\{x_1, \ldots, x_n, x\}]$.
    Since $x$ is not a linear combination of $x_1, \ldots, x_n$ we have $M \subsetneq N$.
    By Exercise 3 part (ii) we have $N \cap M^\perp \ne \{0\}$.
    Let $y \in N \cap M^\perp$ be nonzero. Then $y \in N$ and $y \perp x_i$ for all $i$, and since $y \notin M$ we have
    \[[\{x_1, \ldots, x_n, y\}] = [\{x_1, \ldots, x_n, x\}] = N.\]
\end{proof}
\begin{proof}
    Suppose $x_1, \ldots, x_n$ are linearly independent vectors in $E$.
    We construct $y_1, \ldots, y_n$ inductively.
    Set $y_1 = x_1 / \|x_1\|$, so $[\{y_1\}] = [\{x_1\}]$.
    Given orthonormal $y_1, \ldots, y_k$ with $[\{y_1, \ldots, y_k\}] = [\{x_1, \ldots, x_k\}]$,
    apply part (i) with $x_1, \ldots, x_k$ replaced by $y_1, \ldots, y_k$ and $x$ replaced by $x_{k+1}$
    to obtain a nonzero $z \in [\{y_1, \ldots, y_k, x_{k+1}\}] \cap [\{y_1, \ldots, y_k\}]^\perp$.
    Set $y_{k+1} = z / \|z\|$. Then $y_1, \ldots, y_{k+1}$ are orthonormal and
    \[[\{y_1, \ldots, y_{k+1}\}] = [\{y_1, \ldots, y_k, x_{k+1}\}] = [\{x_1, \ldots, x_{k+1}\}].\]
    By induction, the construction gives orthonormal $y_1, \ldots, y_n$ with $[\{y_1, \ldots, y_k\}] = [\{x_1, \ldots, x_k\}]$ for $k = 1, \ldots, n$.
\end{proof}

\begin{tcolorbox}[title=Problem 9, breakable]
Given a system of homogeneous linear equations
\begin{align*}
    a_{11}x_1 + a_{12}x_2 + \cdots + a_{1n}x_n &= 0 \\
    a_{21}x_1 + a_{22}x_2 + \cdots + a_{2n}x_n &= 0 \\
    &\vdots \\
    a_{m1}x_1 + a_{m2}x_2 + \cdots + a_{mn}x_n &= 0
\end{align*}
with real coefficients $a_{ij}$. If $m < n$, we know from Theorem 3.7.1 that there exist nonzero solution-vectors $x = (x_1, \ldots, x_n)$; deduce another proof of this from Corollary 5.2.8. \{Hint: Let $A = (a_{ij})$ be the $m \times n$ matrix of coefficients and let $M$ be the row space of $A$ (4.2.9).\}
\end{tcolorbox}

\begin{proof}
    Suppose $m < n$.
    Let $A = (a_{ij})$ be the $m \times n$ matrix of coefficients and let $M$ be the row space of $A$, a subspace of $\mathbb{R}^n$.
    By Corollary 5.2.8 we have
    \[\dim M^\perp = n - \dim M.\]
    Since $M$ is spanned by $m$ row vectors, $\dim M \leq m < n$, thus
    \[\dim M^\perp = n - \dim M \geq n - m > 0.\]
    So $M^\perp \ne \{0\}$; let $x \in M^\perp$ be nonzero.
    Then $x$ is orthogonal to every row of $A$ so $Ax = 0$.
\end{proof}


\begin{tcolorbox}[title=Problem 10, breakable]
A \textit{unitary space} is a finite-dimensional complex vector space $E$ with an `inner product' $(x, y) \mapsto (x \mid y)$ taking complex values, satisfying (1) and (3) of Definition 5.1.1, (4) for complex scalars, and, instead of (2), the identity $(x \mid y) = \overline{(y \mid x)}$, where the overbar denotes complex-conjugate (cf.\ \S5.1, Exercise 11).
\begin{enumerate}[(i)]
    \item Show that $(x \mid cy) = \bar{c}(x \mid y)$ for all vectors $x$, $y$ and all complex numbers $c$.
    \item If $y \in E$ define $y'(x) = (x \mid y)$ for all $x \in E$. Show that $y' \in E' = \mathcal{L}(E, \mathbf{C})$ and that $(y + z)' = y' + z'$, $(cy)' = \bar{c}y'$.
    \item Prove the unitary space analogue of Theorem 5.2.5, with `isomorphism' replaced by `conjugate-isomorphism' in the sense suggested by (ii).
    \item Prove the unitary space analogues of the remaining results in this section.
    \item Verify the `Polarization identity'
    \[
        (x \mid y) = \tfrac{1}{4}\{\|x + y\|^2 - \|x - y\|^2 + i\|x + iy\|^2 - i\|x - iy\|^2\}.
    \]
\end{enumerate}
\end{tcolorbox}

\begin{proof}
    We have
    \[\bar{c}(x \mid y) 
        = \bar{c}\,\overline{(y \mid x)} 
        = \overline{c(y \mid x)} 
        = \overline{(cy \mid x)} = (x \mid cy).\]
\end{proof}

\begin{proof}
    Suppose $y \in E$ and define $y'(x) = (x \mid y)$ for all $x \in E$. 
    Additivity still holds so
    Then
    \[(y+z)'(x) = (x \mid y+z) = (x \mid y) + (x \mid z) = y'(x) + z'(x).\]
    Then
    \[(cy)'(x) = (x \mid cy) = \bar{c}(x \mid y) = \bar{c}\,y'(x).\]
\end{proof}

\begin{proof}
    The map $y \mapsto y'$ is linear by (ii). 
    Suppose $y' = 0$ then $(x \mid y) = 0$ for all $x \in E$, 
        so taking $x = y$ gives $(y \mid y) = 0$, thus $y = 0$. 
    Since $\dim E' = \dim E$ (3.9.2), the injective map 
        $y \mapsto y'$ is bijective, thus an isomorphism $E \to E'$.
\end{proof}

\textbf{Solution (iv): Skip, very boring.}

\begin{proof}
    We use $(x \mid x) = \|x\|^2$:
    \begin{align*}
        \|x+y\|^2 &= \|x\|^2 + (x\mid y) + (y \mid x) + \|y\|^2,\\
        \|x-y\|^2 &= \|x\|^2 - (x\mid y) - (y \mid x) + \|y\|^2,\\
        \|x+iy\|^2 &= \|x\|^2 + (x\mid iy) + (iy\mid x) + \|y\|^2 = \|x\|^2 - i(x\mid y) + i(y\mid x) + \|y\|^2,\\
        \|x-iy\|^2 &= \|x\|^2 + i(x\mid y) - i(y\mid x) + \|y\|^2.
    \end{align*}
    Thus
    \begin{align*}
        \|x+y\|^2 - \|x-y\|^2 &= 2(x\mid y) + 2(y\mid x),\\
        i\|x+iy\|^2 - i\|x-iy\|^2 &= 2(x\mid y) - 2(y\mid x).
    \end{align*}
    Summing and multiplying by $\tfrac{1}{4}$ gives $(x \mid y)$.
\end{proof}


\begin{tcolorbox}[title=Problem 11, breakable]
If $E$ is a unitary space (Exercise 10) and $T : E \to E$ is a linear mapping such that $(Tx \mid x) = 0$ for all vectors $x$, then (in contrast with the example in 5.2.4) $T = 0$. \{Hint: Adapt the argument for (v) of Exercise 10 to express $(Tx \mid y)$ as a linear combination of numbers $(Tz \mid z)$.\}
\end{tcolorbox}

\begin{proof}
    Suppose $E$ is a unitary space (Exercise 10)
        and $T : E \to E$ is a linear mapping
        such that $(Tx \mid x) = 0$ for all vectors $x$.
    For $x, y \in E$
    \[(Tx \mid y) = \tfrac{1}{4}\{(T(x+y)\mid x+y) - (T(x-y)\mid x-y) + i(T(x+iy)\mid x+iy) - i(T(x-iy)\mid x-iy)\},\]
    Since $(Tz \mid z) = 0$ for all $z$ we have $(Tx \mid y) = 0$ for all $x, y \in E$. 
    Suppose $y = Tx$ then $(Tx \mid Tx) = \|Tx\|^2 = 0$.
    Thus $Tx = 0$ for all $x$, so $T = 0$.
\end{proof}

\begin{tcolorbox}[title=Problem 12, breakable]
Let $E$ be a Euclidean space.
\begin{enumerate}[(i)]
    \item If $u$, $x$ are vectors in $E$ such that $\|u\| = 1$ and $x$ is not a multiple of $u$, find a nonzero vector $z$ such that $u \perp z$ and $x = au + z$ for a suitable scalar $a$.

    \item If $x_1$, $x_2$ are linearly independent vectors, find orthonormal vectors $u_1$, $u_2$ such that $u_1$ is a multiple of $x_1$, and $u_2$ is a linear combination of $x_1$, $x_2$.

    \item If $x_1, \ldots, x_n$ are linearly independent, then there exist orthonormal vectors $u_1, \ldots, u_n$ such that $[\{x_1, \ldots, x_k\}] = [\{u_1, \ldots, u_k\}]$ for $k = 1, \ldots, n$. (\textit{Gram--Schmidt orthogonalization process})

    \item Infer an alternate proof of 5.2.16: Every Euclidean space has an orthonormal basis.
\end{enumerate}
\{Hints:
\begin{enumerate}[(i)]
    \item If such $a$, $z$ exist, then $(x \mid u) = a$ and $z = x - au$.
    \item Let $u_1 = \|x_1\|^{-1}x_1$, apply (i) with $u = u_1$, $x = x_2$ and normalize the resulting vector $z$.
    \item Induction.\}
\end{enumerate}
\end{tcolorbox}

\begin{proof}
    Set $a = (x \mid u)$ and $z = x - au$. 
    Then $z \neq 0$ since $x$ is not a multiple of $u$, 
        and $(u \mid z) = (u \mid x) - a(u \mid u) = a - a = 0$.
\end{proof}

\begin{proof}
    Set $u_1 = \|x_1\|^{-1}x_1$. 
    Apply (i) with $u = u_1$, $x = x_2$ 
        to get nonzero $z = x_2 - (x_2 \mid u_1)u_1$ with $u_1 \perp z$. 
    Set $u_2 = \|z\|^{-1}z$.
\end{proof}

\begin{proof}
    By induction. 
    The base case $k=1$ is $u_1 = \|x_1\|^{-1}x_1$. 
    Given orthonormal $u_1,\ldots,u_k$ 
        with $[\{u_1,\ldots,u_k\}]=[\{x_1,\ldots,x_k\}]$, 
        set $z = x_{k+1} - \sum_{i=1}^k (x_{k+1}\mid u_i)u_i$ 
        and $u_{k+1} = \|z\|^{-1}z$.
\end{proof}

\begin{proof}
    Every Euclidean space has a 
        (finite) basis of linearly independent vectors, then apply (iii).
\end{proof}

\begin{tcolorbox}[title=Problem 13, breakable]
If $u_1, \ldots, u_k$ are orthonormal vectors in a Euclidean space $E$ of dimension $n > k$, then $E$ has an orthonormal basis $u_1, \ldots, u_k, u_{k+1}, \ldots, u_n$. \{Hint: Consider $M^\perp$ for the appropriate linear subspace $M$.\}
\end{tcolorbox}

\begin{proof}
    Suppose $u_1, \ldots, u_k$
        are orthonormal vectors in a Euclidean space $E$ of dimension $n > k$.
    Since $u_1, \ldots, u_k$ are linearly independent,
        we can expand the list to a basis $u_1, \ldots, u_k, a_{k+1}, \ldots, a_n$ of $E$.
    Applying Problem 8 to $a_{k+1}, \ldots, a_n$ with respect to the already orthonormal $u_1, \ldots, u_k$ gives orthonormal $u_{k+1}, \ldots, u_n$.
    thus $u_1, \ldots, u_k, u_{k+1}, \ldots, u_n$ is an orthonormal basis of $E$.
\end{proof}

\begin{tcolorbox}[title=Problem 14, breakable]
If $M$ is a linear subspace of a 
Euclidean space $E$ then $\dim M + \dim M^\perp = \dim E$ 
by the `alternate proof' of Corollary 5.2.8. 
Infer from $M \cap M^\perp = \{\theta\}$ that $E = M + M^\perp$. \{Hint: \S3.5, Exercise 26 or \S3.6, Exercise 5.\}
\end{tcolorbox}

\begin{proof}
    Suppose $M$ is a linear subspace of a Euclidean space $E$.
    We know that $M \cap M^\perp = \{\theta\}$.
    Since $M, M^\perp$ are linear subspaces of $E$ and $\dim M + \dim M^\perp = \dim E$, 
    it follows directly from Section 3.5 Exercise 26 that $E = M + M^\perp$.
\end{proof}

\subsection{The Adjoint of a Linear Mapping}

\begin{tcolorbox}[title=Problem 2, breakable]
    If $E$ and $F$ are Euclidean spaces, 
    then the mapping $T \mapsto T^*$ 
    is a linear bijection $\mathcal{L}(E,F) \to \mathcal{L}(F,E)$. 
    What is the inverse mapping?
\end{tcolorbox}

\begin{proof}
    Suppose $E$ and $F$ are Euclidean spaces. Consider the diagram below
    \[
        \begin{tikzcd}
        E \arrow[d, "J_E"'] & F \arrow[l, "T^*"'] \arrow[d, "J_F"] \\
        E' & F' \arrow[l, "T'"]
        \end{tikzcd}
    \]
    Since $J_E$, $T'$, $J_F$ are linear isomorphisms, 
    $T^* = J_E^{-1} T' J_F$ is a linear bijection. 
    The inverse mapping is $J^{-1} _FT'^{-1} J_E$.
\end{proof}

\begin{proof}
    By Theorem 5.3.3, $(T^*)^* = T$, so the mapping $T \mapsto T^*$ is its own inverse, thus a bijection. Linearity follows from Theorem 5.3.3 as well. Thus $T \mapsto T^*$ is a linear bijection with inverse $S \mapsto S^*$.
\end{proof}

\begin{tcolorbox}[title=Problem 3, breakable]
    If $E$ and $F$ are inner product spaces 
        and if $T:E \to F$, $S:F \to E$ 
        are mappings such that $(Tx \mid y) = (x \mid Sy)$ 
        for all $x \in E$ and $y \in F$, 
        then $S$ and $T$ are linear
         (and $S = T^*$ when $E$ and $F$ are Euclidean spaces).
\end{tcolorbox}

\begin{proof}
    Suppose $E$ and $F$ are inner product spaces
    and $T:E \to F$, $S:F \to E$
    are mappings such that $(Tx \mid y) = (x \mid Sy)$
    for all $x \in E$ and $y \in F$.
    Let $x_1, x_2 \in E$ and $y \in F$.
    Then
    \[(T(x_1 + x_2) \mid y) = (x_1 + x_2 \mid Sy) = (x_1 \mid Sy) + (x_2 \mid Sy).\]
    Also
    \[(Tx_1 + Tx_2 \mid y) = (Tx_1 \mid y) + (Tx_2 \mid y) = (x_1 \mid Sy) + (x_2 \mid Sy).\]
    So
    \[(T(x_1 + x_2) - (Tx_1 + Tx_2) \mid y) = 0 \text{ for all } y \in F.\]
    Thus $T(x_1 + x_2) = Tx_1 + Tx_2$. Similarly, $T(cx) = cTx$. Thus $T$ is linear.
    The argument for $S$ is symmetrical. When $E$ and $F$ are Euclidean spaces,
    $(Tx \mid y) = (x \mid Sy)$ is precisely the condition of Theorem 5.3.2 (4) to conclude $S = T^*$.
\end{proof}


\begin{tcolorbox}[title=Problem 4, breakable]
For $A$, $B$ in the vector space $\mathrm{M}_n(\mathbf{R})$ of $n \times n$ real matrices (4.1.5), define
\[
    (A \mid B) = \mathrm{tr}(AB').
\]
Prove that $\mathrm{M}_n(\mathbf{R})$ is a Euclidean space (of dimension $n^2$) and find the formula for $\|A\|$ if $A = (a_{ij})$.
\end{tcolorbox}

\begin{proof}
    We have $\mathrm{tr}(AB') = \mathrm{tr}((AB')') = \mathrm{tr}(BA')$, 
        so $(A\mid B) = (B\mid A)$. 
    Linearity follows from linearity of the trace and transpose. 
    If $A \ne 0$ we have
    \[(A \mid A) = \mathrm{tr}(AA') = \sum_{i,j} a_{ij}^2 \geq 0.\]
    Thus $\mathrm{M}_n(\mathbf{R})$ is a Euclidean space. 
    Suppose $A$ is the matrix over the linear transformation $T : V \longrightarrow W$.
    We have $\dim V = \dim W = n$ so $\dim A = (\dim V)(\dim W) = n^2$. Finally,
    \[\|A\| = \left(\sum_{i,j=1}^n a_{ij}^2\right)^{\frac{1}{2}}.\]
\end{proof}

\begin{tcolorbox}[title=Problem 5, breakable]
The converse of Theorem 5.3.6 is also true, thus $T(M) \subset N$ if and only if $T^*(N^\perp) \subset M^\perp$.
\end{tcolorbox}

\begin{proof}
    Let $E$ and $F$ be Euclidean spaces, $T: E \to F$ a linear mapping,
    $M$ a linear subspace of $E$, and $N$ a linear subspace of $F$.
    Furthermore, suppose $T^*(N^\perp) \subset M^\perp$.
    Let $x \in M$. We must show $Tx \in N$.
    It suffices to show $(Tx \mid y) = 0$ for all $y \in N^\perp$.
    For any $y \in N^\perp$, we have $T^*y \in M^\perp$ by hypothesis, so
    \[(Tx \mid y) = (x \mid T^*y) = 0,\]
    since $x \in M$ and $T^*y \in M^\perp$. Thus $Tx \in N^{\perp\perp} = N$.
\end{proof}


\begin{tcolorbox}[title=Problem 6, breakable]
Theorem 5.3.6 looks suspiciously like \S3.9, Exercise 18 in disguise. Can you unmask it?
\end{tcolorbox}

\begin{theorem}[Theorem 5.3.6]
    Let $E$ and $F$ be Euclidean spaces, $T : E \longrightarrow F$
    a linear mapping, $M$ a linear subspace of $E$, and $N$ a linear subspace 
    of $F$. 

    If $T(M) \subset N$ then $T^*(N^\perp) \subset M^\perp$.
\end{theorem}

\begin{theorem}[Exercise 18]
    Let $T : V \longrightarrow W$ be a linear mapping, $M$ a linear subspace 
    of $V$, $N$ a linear subspace of $W$. 

    If $T(M) \subset N$ then $T'(N^\circ) \subset M^\circ$.
\end{theorem}

\begin{proof}
    The annihilator 
    $N^\circ = \{ f \in W' : f(n) = 0 \ \forall\, n \in N \}$ corresponds to 
    $N^\perp = \{ v \in F : (v \mid n) = 0 \ \forall\, n \in N \}$.
    $T'$ takes a linear form on $W$ and produces a linear form on $V$,
    whereas $T^*$ takes a vector in $F$ and produces a vector in $E$,
    with the inner product $(w \mid \cdot)$ playing the role of the linear form.
\end{proof}


\begin{tcolorbox}[title=Problem 7, breakable]
Let $E$ and $F$ be Euclidean spaces, $T:E \to F$ a linear mapping, $A$ a subset of $E$, and $B$ a subset of $F$. Prove:
\begin{enumerate}[(i)]
    \item $(T(A))^\perp = (T^*)^{-1}(A^\perp)$;
    \item[(i')] $(T^*(B))^\perp = T^{-1}(B^\perp)$;
    \item[(ii)] $(T(E))^\perp = \ker T^*$, therefore $F = T(E) \oplus \ker T^*$;
    \item[(ii')] $(T^*(F))^\perp = \ker T$, therefore $E = T^*(F) \oplus \ker T$;
    \item[(iii)] $T(A) \subset B \Rightarrow T^*(B^\perp) \subset A^\perp$.
\end{enumerate}
\end{tcolorbox}

\begin{proof}
\textbf{(i)} $v \in (T(A))^\perp \iff (v \mid T(a)) = 0\ \forall a \in A 
\iff (T^*v \mid a) = 0\ \forall a \in A \iff T^*v \in A^\perp 
\iff v \in (T^*)^{-1}(A^\perp)$.

\textbf{(i')} $v \in (T^*(B))^\perp \iff (v \mid T^*b) = 0\ \forall b \in B
\iff (Tv \mid b) = 0\ \forall b \in B \iff Tv \in B^\perp
\iff v \in T^{-1}(B^\perp)$.

\textbf{(ii)} Set $A = E$ in (i): $(T(E))^\perp = (T^*)^{-1}(E^\perp) = (T^*)^{-1}(\{0\}) = \ker T^*$.
Since $T(E)$ is a subspace of $F$, we have $F = T(E) \oplus (T(E))^\perp = T(E) \oplus \ker T^*$.

\textbf{(ii')} Set $B = F$ in (i'): $(T^*(F))^\perp = T^{-1}(F^\perp) = T^{-1}(\{0\}) = \ker T$.
Since $T^*(F)$ is a subspace of $E$, we have $E = T^*(F) \oplus (T^*(F))^\perp = T^*(F) \oplus \ker T$.

\textbf{(iii)} Let $v \in B^\perp$ and $a \in A$. Since $T(a) \in T(A) \subset B$,
$(T^*v \mid a) = (v \mid T(a)) = 0$, so $T^*v \in A^\perp$.
\end{proof}

\subsection{Orthogonal Mappings and Matrices}

\begin{tcolorbox}[title=Problem 1, breakable]
Let $E$ be a Euclidean space, $T \in \mathcal{L}(E)$. Show the following conditions are equivalent:
\begin{enumerate}[(a)]
    \item $T$ is orthogonal (that is, $T$ preserves inner products);
    \item $\|Tx - Ty\| = \|x - y\|$ for all vectors $x, y$ (that is, $T$ preserves distances).
\end{enumerate}
\end{tcolorbox}

\begin{proof}
    Suppose $T$ is orthogonal. Let $x, y \in E$. Then
    \[\|Tx - Ty\|^2 = (Tx - Ty \mid Tx - Ty) = (T(x-y) \mid T(x-y)) = (x-y \mid x-y) = \|x-y\|^2.\]
    Thus $\|Tx - Ty\| = \|x - y\|$.

    Conversely, suppose $\|Tx - Ty\| = \|x - y\|$ for all $x, y \in E$. Squaring both sides,
    \[(T(x-y) \mid T(x-y)) = (x-y \mid x-y)\]
    and since $x - y$ ranges over all of $E$, $T$ preserves inner products.
\end{proof}

\begin{tcolorbox}[title=Problem 2, breakable]
Let $x_1, \ldots, x_n$ be an orthonormal basis of a Euclidean space $E$, let $r_1, \ldots, r_n$ be real numbers whose squares are equal, and let $T \in \mathcal{L}(E)$ be the linear mapping such that $Tx_i = r_i x_i$ for all $i$. Prove that $T$ preserves orthogonality, in the sense that $(x \mid y) = 0 \Rightarrow (Tx \mid Ty) = 0$, and that $T$ is a scalar multiple of an orthogonal mapping, but $T$ is not orthogonal unless $r_i = \pm 1$ for all $i$.
\end{tcolorbox}

\begin{proof}
    Let $x = \sum_i a_i x_i$ and $y = \sum_i b_i x_i$ and suppose $(x \mid y) = 0$, so $\sum_i a_i b_i = 0$. Then
    \[(Tx \mid Ty) = \left(\sum_i a_i r_i x_i \,\middle|\, \sum_i b_i r_i x_i\right) = \sum_i a_i b_i r_i^2 = r^2 \sum_i a_i b_i = 0\]
    where $r^2 = r_i^2$ is the common value. Thus $T$ preserves orthogonality.

    Since $r^2$ is common, $T = rU$ where $U x_i = \frac{r_i}{r} x_i$ with $\frac{r_i}{r} = \pm 1$,
    so $U$ is orthogonal and $T$ is a scalar multiple of an orthogonal mapping.

    Finally, $T$ is orthogonal iff $(Tx_i \mid Tx_i) = (x_i \mid x_i) = 1$, thus \ $r_i^2 = 1$,
    so \ $r_i = \pm 1$ for all $i$.
\end{proof}

\newpage
\begin{tcolorbox}[title=Problem 3, breakable]
If $E$ is a Euclidean space and $T \in \mathcal{L}(E)$ 
is a linear mapping that preserves orthogonality (in the sense of Problem 2), 
prove that $T$ is a scalar multiple of an orthogonal mapping:
\begin{enumerate}[(i)]
    \item If $\|x\| = \|y\|$ then $\|Tx\| = \|Ty\|$.
    \item Let $x_1, \ldots, x_n$ be an orthonormal basis of $E$ 
    and let $r$ be the common value of $\|Tx_i\|$. Show that $T^*T = r^2 I$.
    \item If $r \neq 0$ then $r^{-1}T$ is orthogonal.
\end{enumerate}
\end{tcolorbox}

\begin{proof}
    Suppose $\|x\| = \|y\|$.
    Then $x + y \perp x - y$ since
    \[(x+y \mid x-y) = \|x\|^2 - \|y\|^2 = 0.\]
    Since $T$ preserves orthogonality, $T(x+y) \perp T(x-y)$. Then
    \[0 = (T(x+y) \mid T(x-y)) = \|Tx\|^2 - \|Ty\|^2,\]
    so $\|Tx\| = \|Ty\|$.
\end{proof}

\begin{proof}
    For all $i, j$ we have
    \[(T^* Tx_i \mid x_j) = (Tx_i \mid Tx_j).\]
    Sifnce $x_i \perp x_j$ for $i \ne j$, $T$ preserves orthogonality gives $Tx_i \perp Tx_j$,
    so $(Tx_i \mid Tx_j) = 0$ for $i \ne j$.
    For $i = j$, $(Tx_i \mid Tx_i) = \|Tx_i\|^2 = r^2$.
    Thus $(T^*T x_i \mid x_j) = r^2 \delta_{ij} = (r^2 I x_i \mid x_j)$
    for all $i, j$, so $T^*T = r^2 I$.
\end{proof}

\begin{proof}
    Suppose $r \ne 0$.
    Let $x, y \in E$. Then
    \[(r^{-1} Tx \mid r^{-1} Ty) = r^{-2}(Tx \mid Ty) = r^{-2}(T^*Tx \mid y) = r^{-2} \cdot r^2(x \mid y) = (x \mid y).\]
    Thus $r^{-1} T$ is orthogonal.
\end{proof}

\begin{tcolorbox}[title=Problem 4, breakable]
Let $A$ and $B$ be $n \times n$ real matrices, $I$ the $n \times n$ identity matrix. Prove:
\begin{enumerate}[(i)]
    \item $I$ is an orthogonal matrix;
    \item if $A$ and $B$ are orthogonal, then so is $AB$;
    \item if $A$ is orthogonal then it is invertible and $A^{-1}$ is orthogonal.
\end{enumerate}
\end{tcolorbox}

\begin{proof}
    Clearly $I^T I = I$, so $I$ is orthogonal.
    
    Suppose $A$ and $B$ are orthogonal, so $A^T A = I$ and $B^T B = I$. Then
    \[(AB)^T(AB) = B^T A^T A B = B^T I B = B^T B = I,\]
    so $AB$ is orthogonal.

    Suppose $A$ is orthogonal, so $A^T A = I$, meaning $A^{-1} = A^T$. Then
    \[(A^{-1})^T A^{-1} = (A^T)^T A^T = A A^T = I,\]
    so $A^{-1}$ is orthogonal.
\end{proof}

\begin{tcolorbox}[title=Problem 5, breakable]
State and prove the analogue of Problem 4 for linear mappings.
\end{tcolorbox}

\begin{proof}
    Clearly $I'I = I$, so $I$ is orthogonal.

    Suppose $S$ and $T$ are orthogonal, so $T'T = I$ and $S'S = I$. Then
    \[(ST)'(ST) = T' S' S T = T' I T = T' T = I,\]
    so $ST$ is orthogonal.

    Suppose $T$ is orthogonal, so $T'T = I$, meaning $T^{-1} = T'$. Then
    \[(T^{-1})' T^{-1} = (T')' T' = T T' = I,\]
    so $T^{-1}$ is orthogonal.
\end{proof}

\newpage
\begin{tcolorbox}[title=Problem 6, breakable]
If $T \in \mathcal{L}(\mathbb{R}^2)$ is orthogonal, show that $T$ may be interpreted as either a rotation or a reflection in the Euclidean plane.
\end{tcolorbox}

\begin{proof}
    Suppose $T \in \mathcal{L}(\mathbb{R}^2)$ is orthogonal.
    Let $A$ be the matrix of $T$ and suppose 
    \[A = \begin{pmatrix}
        a & b \\ c & d
    \end{pmatrix}.\]
    Since $A^T A = I$ we have 
    \[\begin{pmatrix}a&b\\c&d\end{pmatrix}^T\begin{pmatrix}a&b\\c&d\end{pmatrix} 
    = \begin{pmatrix}a^2+c^2 & ab+cd \\ ab+cd & b^2+d^2\end{pmatrix} = I\]
    So $a^2 + c^2 = 1$, $ab + cd = 0$, $b^2 + d^2 = 1$.
    From $a^2 + c^2 = 1$ so $a = \cos\theta$, $c = \sin\theta$ for some $\theta$,
    and combined with $b^2 + d^2 = 1$ we have two cases.

    \textit{Case 1:} $b = -\sin\theta$, $d = \cos\theta$. Then
    \[A = \begin{pmatrix}\cos\theta & -\sin\theta \\ \sin\theta & \cos\theta\end{pmatrix},\]
    which is a rotation by $\theta$.

    \textit{Case 2:} $b = \sin\theta$, $d = -\cos\theta$. Then
    \[A = \begin{pmatrix}\cos\theta & \sin\theta \\ \sin\theta & -\cos\theta\end{pmatrix},\]
    which is a reflection through the line at angle $\theta/2$.

    Conversely, for any rotation by $\theta$ or reflection through a line at angle $\theta/2$
    so
    \[\begin{pmatrix}\cos\theta & \sin\theta \\ -\sin\theta & \cos\theta\end{pmatrix}
    \begin{pmatrix}\cos\theta & -\sin\theta \\ \sin\theta & \cos\theta\end{pmatrix}
    = \begin{pmatrix}1 & 0 \\ 0 & 1\end{pmatrix} = I,\]
    \[\begin{pmatrix}\cos\theta & \sin\theta \\ \sin\theta & -\cos\theta\end{pmatrix}
    \begin{pmatrix}\cos\theta & \sin\theta \\ \sin\theta & -\cos\theta\end{pmatrix}
    = \begin{pmatrix}1 & 0 \\ 0 & 1\end{pmatrix} = I.\]
    Thus every rotation and every reflection is orthogonal.
\end{proof}

\begin{tcolorbox}[title=Problem 7, breakable]
Let $E$ be a Euclidean space, $T \in \mathcal{L}(E)$ an orthogonal linear mapping, and $M$ a linear subspace of $E$ such that $T(M) \subset M$. Prove that $T(M) = M$ and $T(M^\perp) = M^\perp$.
\end{tcolorbox}

\begin{proof}
    Since $T$ is an orthogonal linear mapping, $\dim T(M) = \dim M$.
    Since $T(M) \subset M$ we have $T(M) = M$.

    Now let $x \in M^\perp$ and $y \in M$. Since $T(M) = M$ there exists $m \in M$ such that $y = Tm$, and so
    \[(Tx \mid y) = (Tx \mid Tm) = (x \mid m) = 0,\]
    where we used that $T$ is orthogonal. Thus $Tx \in M^\perp$, so $T(M^\perp) \subset M^\perp$.
    Applying the same argument to $M^\perp$ gives $T(M^\perp) = M^\perp$.
\end{proof}

\begin{tcolorbox}[title=Problem 8, breakable]
Show that the matrix
\[\frac{1}{3}\begin{pmatrix} 1 & -2 & 2 \\ 2 & 2 & 1 \\ 2 & -1 & -2 \end{pmatrix}\]
is orthogonal.
\end{tcolorbox}

\[\frac{1}{9}\begin{pmatrix}1&-2&2\\2&2&1\\2&-1&-2\end{pmatrix}\begin{pmatrix}1&2&2\\-2&2&-1\\2&1&-2\end{pmatrix} = \frac{1}{9}\cdot 9I = I\]

\newpage
\begin{tcolorbox}[title=Problem 12, breakable]
Let $e_1, e_2$ be the canonical basis of $E^2$, $x_1 = e_1 + e_2$, $x_2 = e_1 - e_2$, and let $A = \{x_1, x_2, -x_1, -x_2\}$; interpret $A$ as the set of vertices of a square, centered at the origin, with edge length 2. Let $O = \{T \in \mathcal{L}(E^2) : T(A) = A\}$, the set of all linear mappings of Euclidean 2-space that preserve the vertices of the square. Prove:
\begin{enumerate}[(i)]
    \item Every $T \in O$ is bijective.
    \item $O$ contains 8 elements.
    \item $I \in O$; if $S, T \in O$ then $ST \in O$ and $T^{-1} \in O$.
    \item Every $T \in O$ is orthogonal.
    \item Find the matrices of the $T \in O$ relative to the canonical basis; relative to the basis $x_1, x_2$.
    \item Interpret the elements of $O$ as $I, R, R^2, R^3, H, V, D, E$, where $R$ is counter-clockwise rotation of $90^\circ$, $H$ is reflection in the $x$-axis, $V$ is reflection in the $y$-axis, $D$ is reflection in the diagonal containing $\pm x_1$, and $E$ is reflection in the other diagonal.
\end{enumerate}
\end{tcolorbox}

\begin{proof}
    Since $\dim T(A) = \dim A$, $T$ is surjective and thus bijective.
\end{proof}

\begin{proof}
    $T$ is determined by where it sends $x_1$ and $x_2$. We need $T(x_1) \in A$ and $T(x_2) \in A$ with $\{T(x_1), T(\-x_1), T(x_2), T(-x_2)\} = A$. There are 4 choices for $T(x_1)$ and then 2 remaining choices for $T(x_2)$, giving $4 \times 2 = 8$ elements.
\end{proof}

\begin{proof}
    Clearly $I(A) = A$ so $I \in O$. 
    Suppose $S, T \in O$
    Then $ST(A) = S(T(A)) = S(A) = A$, so $ST \in O$. 
    Since $T$ is bijective, $T^{-1}(A) = A$, so $T^{-1} \in O$.
\end{proof}

\begin{proof}
    Every $T \in O$ permutes $\{x_1, x_2, -x_1, -x_2\}$, 
    so it maps $\{x_1, -x_1\}$ and $\{x_2, -x_2\}$ 
    to pairs of orthogonal vertices. 
    We have $\|Tx_i\| = \|x_i\| = \sqrt{2}$ 
        and $Tx_1 \perp Tx_2$. 
    Thus $T$ preserves norms and orthogonality, thus is orthogonal.
\end{proof}

\begin{proof}
    Relative to the canonical basis $e_1, e_2$, using $x_1 = e_1 + e_2$, $x_2 = e_1 - e_2$
    \[
    I = \begin{pmatrix} 1 & 0 \\ 0 & 1 \end{pmatrix}, \quad
    R = \begin{pmatrix} 0 & -1 \\ 1 & 0 \end{pmatrix}, \quad
    R^2 = \begin{pmatrix} -1 & 0 \\ 0 & -1 \end{pmatrix}, \quad
    R^3 = \begin{pmatrix} 0 & 1 \\ -1 & 0 \end{pmatrix},
    \]
    \[
    H = \begin{pmatrix} 1 & 0 \\ 0 & -1 \end{pmatrix}, \quad
    V = \begin{pmatrix} -1 & 0 \\ 0 & 1 \end{pmatrix}, \quad
    D = \begin{pmatrix} 0 & 1 \\ 1 & 0 \end{pmatrix}, \quad
    E = \begin{pmatrix} 0 & -1 \\ -1 & 0 \end{pmatrix}.
    \]
    Relative to the basis $x_1, x_2$
    \[
    I = \begin{pmatrix} 1 & 0 \\ 0 & 1 \end{pmatrix}, \quad
    R = \begin{pmatrix} 0 & -1 \\ 1 & 0 \end{pmatrix}, \quad
    R^2 = \begin{pmatrix} -1 & 0 \\ 0 & -1 \end{pmatrix}, \quad
    R^3 = \begin{pmatrix} 0 & 1 \\ -1 & 0 \end{pmatrix},
    \]
    \[
    H = \begin{pmatrix} 0 & 1 \\ 1 & 0 \end{pmatrix}, \quad
    V = \begin{pmatrix} 0 & -1 \\ -1 & 0 \end{pmatrix}, \quad
    D = \begin{pmatrix} 1 & 0 \\ 0 & -1 \end{pmatrix}, \quad
    E = \begin{pmatrix} -1 & 0 \\ 0 & 1 \end{pmatrix}.
    \]
\end{proof}